\section*{Sheet 1}

\section*{Sheet 2}

\subsection*{6. Unpolarized scattering of fermions in the Yukawa theory}

\begin{center}
    \textit{Compute the unpolarized (i.e. averaged over initial and summed over final polarization states) squared tree-level scattering amplitude} \(\vert \mathcal{M} \vert^2\) \textit{unpolarized for elastic scattering of two identical massless Dirac fermions (e.g.} \(e^- e^- \to e^- e^-\) \textit{for} \(m_e \to 0\)\textit{) in the Yukawa theory, as a functions of the Mandelstam invariants and the mass} \(M\) \textit{of the scalar.}
\end{center}

For the Yukawa theory, the interaction Lagrangian is given by
\[
    \mathcal{L}_{\text{int}} = -g \phi \bar{\psi} \psi,
\]
and the vertices are represented by a factor of \(-ig\). The tree-level scattering amplitude for the process
\[
    e^-(p_1) e^-(p_2) \to e^-(p_3) e^-(p_4),
\]
can be computed using the Feynman rules for the Yukawa theory, which involve the exchange of a scalar particle in the \(t\)-channel and \(u\)-channel:
\[
    i \mathcal{M} = \text{draw t-chan + u-chan},
\]
where using Feynman rules, we have
\[
    i \mathcal{M} = (-ig)^2 \left[ \bar{u}_3 u_1 \frac{i}{t - M^2} \bar{u}_4 u_2 - \bar{u}_4 u_1 \frac{i}{u - M^2} \bar{u}_3 u_2 \right] = \mathcal{M}_1 + \mathcal{M}_2,
\]
with \(t = (p_1 - p_3)^2\) and \(u = (p_1 - p_4)^2\), and the minus sign arises due to the fact that the u-channel is obtained by exchanging the two identical fermions in the final state w.r.t. the t-channel, which introduces a minus sign (see extra rules for fermions).

    [...]

\[
    \vert \mathcal{M}_{\text{unpol.}} \vert^2 = \dots
\]

[...]

\[
    \begin{aligned}
        \sum_{s} \dots & = \dots \\
        \sum_{s} \dots & = \dots \\
        \sum_{s} \dots & = \dots \\
        \sum_{s} \dots & = \dots
    \end{aligned}
\]
and the final, more difficult, is given by
\[
    \begin{aligned}
        \sum_{s} \dots & = \dots \\
                       & = \dots
    \end{aligned}
\]

If we now use the limit for \(m_e \to 0\) we can compte:
\[
    \dots
\]
and now, using the properties of the Mandelstam invariants, we can write the final result as a function of \(s\) and \(t\), since \(s + t + u = 0\). Thus we have
\[
    \vert \mathcal{M}_{\text{unpol.}} \vert^2 = \dots
\]
[...]
\[
    \vert \mathcal{M}_{\text{unpol.}} \vert^2 = \dots
\]

[...]

\begin{remark}
    When the total mass approaches zero, i.e. \(M \to 0\), we get a constant matrix element for the unpolarized scattering amplitude, which is given by \(\vert \mathcal{M}_{\text{unpol.}} \vert^2 \to 3 g^4\). This is a consequence of the fact that the propagator of the scalar particle becomes massless, so that it behaves as \(1/t\) or \(1/u\), and since \(t\) and \(u\) are negative, we get a constant contribution to the scattering amplitude.
\end{remark}

\subsection*{7. Decay of a scalar into a fermion-antifermion pair}

\begin{center}
    \textit{In the Yukawa theory, compute the total decay rate} \( \Gamma \) \textit{of a scalar into a fermion antifermion pair, as a function of the masses} \(M\) \textit{of the scalar and} \(m\) \textit{of the Dirac fermion. Which condition must} \(M\) \textit{and} \(m\) \textit{satisfy for this decay to be possible?}
\end{center}

The vertex of this process is given by
\[
    i \mathcal{M} = \text{draw vertex} = -(ig) \bar{u}_1 v_2,
\]
[...]
\[
    \vert \mathcal{M}_{\text{unpol.}} \vert^2 = \dots
\]

Now, computing the square of the momenta of the final state particles, we have
\[
    M^2 = P^2 = (p_1 + p_2)^2 = 2m^2 + 2 p_1 \cdot p_2, \implies 2 p_1 \cdot p_2 = M^2 - 2m^2,
\]
which let us write the unpolarized squared amplitude as
\[
    \vert \mathcal{M}_{\text{unpol.}} \vert^2 = 2 g^2 (M^2 - 4 m^2).
\]

Now using the formula for the differential decay rate of a particle of mass \(M\) into two particles of mass \(m\) we have
\[
    \mathrm{d} \Gamma = \frac{1}{2M} \vert \mathcal{M}_{\text{unpol.}} \vert^2 \mathrm{d} \Pi_2,
\]
where \(\mathrm{d} \Pi_2\) is the two-particle phase space measure, which can be computed keeping in ming that the decay is computed in the center of mass frame of the particle, so that it is given by the expression we explicitly computed in the lectures\TODO{add reference to the lectures}:
\[
    \mathrm{d} \Pi_2 = \mathrm{d} \Omega \frac{1}{16 \pi^2} \frac{\vert \vec{p}_f \vert}{E_{c.m.}}, \quad \begin{dcases}
        E_{c.m.} = M,                                                 \\
        \vert \vec{p}_f \vert = \vert p_1 \vert = \sqrt{E_1^2 - m^2}, \\
        P = \begin{pmatrix}
                M \\
                \vec{0}
            \end{pmatrix}, \quad p_1 = \begin{pmatrix}
                                           E_1 \\
                                           \vec{p}_f
                                       \end{pmatrix}.
    \end{dcases}
\]

We thus have to compute expressions for \(E_1\) and \(\vert \vec{p}_f \vert\) as a function of \(M\) and \(m\). We can do this by using the fact that the momenta of the final state particles are on-shell, so that we have
\[
    \dots
\]
and
\[
    \dots
\]

With these expressions, we can write the total differential decay rate as
\[
    \frac{\mathrm{d} \Gamma}{\mathrm{d} \Omega} = \dots
\]
and since the integral over the solid angle is given by \(\int \mathrm{d} \Omega = 4 \pi\), we can write the total decay rate as
\[
    \Gamma = \int \frac{\mathrm{d} \Gamma}{\mathrm{d} \Omega} \mathrm{d} \Omega = \dots
\]
[...]
\begin{remark}
    The condition for the decay to be possible is given by the fact that the total energy of the final state particles must be greater than the mass of the initial state particle, so that we must have \(M > 2m\). When we are in the limit for \(M \to 2m\), the decay rate vanishes \(\Gamma \to 0\): this is because the phase space volume for the decay vanishes, since the final state particles are produced at rest, so that the decay is kinematically forbidden.
\end{remark}

\subsection*{8. Compton scattering in QED}

% Long exercise, it is instructive but in the exam we will not find something like this

\begin{center}
    \textit{In QED, compute the unpolarized squared tree-level scattering amplitude for elastic scattering of a massless electron and a photon, as a function of the Mandelstam invariants. Use this result to compute the differential cross section} \(\tfrac{\mathrm{d}\sigma}{\mathrm{d}\Omega}\) \textit{for this process in the centre of mass.}
\end{center}

[...]

\[
    \begin{aligned}
        i \mathcal{M} & = \text{draw sum of 2 diagrams}                                                                                                                       \\
                      & = (-ie)^2 \left[ \bar{u}_3 \gamma^{\nu} \frac{\slashed{p_1} + \slashed{p_2}}{s} \gamma^{\mu} u_1 + \dots \right] = i \mathcal{M}_1 + i \mathcal{M}_2,
    \end{aligned}
\]
[...]
\[
    \vert \mathcal{M}_{\text{unpol.}} \vert^2 = \dots
\]
where now we have to compute the \(A\), \(B\) and \(C\) terms. Starting from \(C\) we have
\[
    \begin{aligned}
        C & = \frac{1}{4} \sum_{s} \dots = \text{product of diagrams}                              \\
          & = \frac{e^4}{2 s u}\dots                                                               \\
          & \times (\epsilon^*_2)_\mu (\epsilon_4^*)_\nu (\epsilon_2)_{\rho} (\epsilon_4)_{\sigma}
    \end{aligned}
\]
where the last term is given by \((\epsilon^*_2)_\mu (\epsilon_4^*)_\nu (\epsilon_2)_{\rho} (\epsilon_4)_{\sigma} = g_{\mu \rho} g_{\nu \sigma}\), enforcing
\[
    C = \dots
\]

If we now use the result computed in exercise 5, in which we demonstrated that
\[
    \gamma^{\mu} \gamma^{\nu} \gamma^{\rho} \gamma^{\sigma} \gamma_{\mu} = - 2 \gamma^{\sigma} \gamma^{\rho} \gamma^{\nu},
\]
it easy to see that this implies
\[
    \implies \dots
\]
Thus we can further simplify the expression for \(C\) as
\[
    C = - \frac{e^4}{2 s u} \Tr[ \dots ]
\]
We need to show that \(\gamma^{\mu} \gamma^{\nu} \gamma^{\rho} \gamma_{\mu} = 4 g_{\nu \rho} \mathbb{I}\).
\begin{proof}
    We can start by defining ...
\end{proof}

Thus now we have all the ingredients to compute the final value for the \(C\)-term, which is given by
\[
    \begin{aligned}
        C & = \dots                                   \\
          & = \dots                                   \\
          & = - \frac{4 e^4}{su} (s + t + u)(-t) = 0,
    \end{aligned}
\]
since we are in the massless limit, so that \(s + t + u = 0\). Thus we have that the interference term between the two diagrams vanishes.

We can now move on to the computation of the \(A\)-term, which is given by
\[
    \begin{aligned}
        A & = \frac{1}{4} \sum_{s} \vert \mathcal{M}_1 \vert^2 = \frac{1}{4} \text{ product of diagrams} \\
          & = \frac{e^4}{4 s^2} \Tr[ \dots ]
    \end{aligned}
\]
where now we need to compute the following product of gamma matrices:
\[
    \gamma^{\mu} \gamma^{\nu} \gamma_\mu = \dots
\]
which will enforce
\[
    \begin{aligned}
        A & = \frac{4 e^4}{4 s^2} \Tr[ \dots ]                                                 \\
          & = \frac{4 e^4}{4 s^2} \Tr[\slashed{p}_2 \slashed{p}_1 \slashed{p}_3 \slashed{p}_2] \\
          & = ...                                                                              \\
          & = ... = -2 e^4 \frac{u}{s},
    \end{aligned}
\]
where we used ...

The \(B\)-term can be computed via crossing symmetry, since it is given by the same expression as \(A\) but ...

    [...]

Thus we can compute the unpolarized squared amplitude as
\[
    \vert \mathcal{M}_{\text{unpol.}} \vert^2 = \dots
\]
Again, we are in the center of mass frame
\[
    E_1 = E_2 = E_3 = E_4 = \frac{E_{cm}}{2} = E,
\]
so that we can write the momenta four-vectors as
\[
    p_1 = \begin{pmatrix}
        E \\
        0 \\
        0 \\
        E
    \end{pmatrix}, \quad p_2 = \begin{pmatrix}
        E \\
        0 \\
        0 \\
        -E
    \end{pmatrix}, \quad p_3 = \begin{pmatrix}
        E             \\
        0             \\
        E \sin \theta \\
        E \cos \theta
    \end{pmatrix}, \quad p_4 = \begin{pmatrix}
        E              \\
        0              \\
        -E \sin \theta \\
        -E \cos \theta
    \end{pmatrix},
\]
and thus the Mandelstam invariants read
\[
    \begin{aligned}
        u & = (p_4 - p_1)^2 = - 2 p_1 \cdot p_4 = -2 (E^2 + E^2 \cos \theta) = -2 E^2 (1 + \cos \theta), \\
        s & = E^2_{cm} = 4 E^2,                                                                          \\
        \frac{u}{s} = - \frac{1 + \cos \theta}{2}.
    \end{aligned}
\]
So that we can write the unpolarized squared amplitude as a function of the scattering angle \(\theta\) as
\[
    \vert \mathcal{M}_{\text{unpol.}} \vert^2 = 2 e^4 \left( \frac{1 + \cos \theta}{2} + \frac{2}{1 + \cos \theta} \right).
\]
Finally, we can compute the differential cross section for this process in the center of mass frame as
\[
    \begin{aligned}
        \frac{\mathrm{d} \sigma}{\mathrm{d} \Omega} & = \dots \\
                                                    & = \dots
    \end{aligned}
\]


\subsection*{9. Ward identity in sQED}

\begin{center}
    \textit{In sQED, explicitly show that the Ward identity is valid for an amplitude of Compton scattering (i.e. the same we did in the lecture but replacing the fermion with a scalar). Note that the amplitude is a sum of three diagrams.}
\end{center}

\[
    \text{Draw of vertex} = \dots
\]
\[
    \text{Draw of vertex} = \dots
\]

\[
    \mathrm{d} \gamma = \phi \gamma = \text{Draw of generic diagram} = \dots
\]

We want to show that \((p_2)_\mu i \mathcal{M}^{\mu \nu} = 0\) at tree level, where \(p_2\) is the momentum of the outgoing photon. We can do this by explicitly computing the amplitude for each diagram and then summing them up, so that we have
\[
    \begin{aligned}
        (p_2)_\mu i \mathcal{M}^{\mu \nu} & = (p_2)_\mu (\text{sum of diagrams}) \\
                                          & = \dots                              \\
                                          & + \dots
    \end{aligned}
\]
where now we need to compute the following terms:
\[
    p_1 \cdot p_2 = \frac{s - m^2}{2}, \quad p_2 \cdot p_4 = -\frac{t}{2}, \quad p_2 \cdot p_3 = - \frac{u - m^2}{2},
\]
and with these expressions we can compute
\[
    \begin{aligned}
        (p_2)_\mu [2p_1^{\mu} + p_2^{\mu}]            & = s - m^2 + 0 = s - m^2,                             \\
        (p_2)_\mu [p_1^{\mu} - p_4^{\mu} + p_3^{\mu}] & = \frac{1}{2} (s - m^2 + t - u + m^2) = - (u - m^2),
    \end{aligned}
\]
where for the last equality we used the fact that \(s + t + u = 2 m^2\). Thus we have
\[
    \begin{aligned}
        (p_2)_\mu i \mathcal{M}^{\mu \nu} & = e^2 \left[ -i (p_1^{\nu} + p_2^{\nu} + p_3^{\nu}) + i (2 p_1^{\nu} - p_4^{\nu}) + 2i p_2^{\nu} \right] \\
                                          & = ie^2 (p_1^{\nu} + p_2^{\nu} - p_3^{\nu} - p_4^{\nu}) = 0,
    \end{aligned}
\]
where the last equality follows from the conservation of momentum at the vertex, which implies that \(p_1 + p_2 = p_3 + p_4\).

\begin{remark}
    All of the three diagrams are necessary to satisfy the Ward identity and thus the cancellation to happen: if we consider only one or two of them, we would get a non-zero contribution to \((p_2)_\mu i \mathcal{M}^{\mu \nu}\), which would violate the Ward identity. The Ward identity is a consequence of gauge invariance, which is a fundamental symmetry of the theory, and thus must be satisfied by all physical amplitudes.
\end{remark}

\subsection*{10. Scaling Green functions and scattering amplitudes}

\begin{center}
    \textit{Show that, under a transformation of the fields of the form}
    \[
        \phi(x) \to c \phi(x),
    \]
    \textit{where} \(c\) \textit{is a constant, then:}
    \begin{enumerate}[label=(\alph*)]
        \item \textit{The Greeen functions of the theory scale as}
              \[
                  \dots
              \]
        \item \textit{The amputated Green functions of the theory scale as}
              \[
                  \dots
              \]
        \item \textit{Scattering matrix elements are invariant under such transformation.}
    \end{enumerate}
    Hint: \textit{Use the LSZ formula and recall that} \(\sqrt{Z} \equiv \bra{\Omega} \phi(0) \ket{\mathbf{p}} \big|_{\mathbf{p} = \mathbf{0}}\).
\end{center}

Defining
\[
    G_N (x_1, \dots, x_N) = \bra{\Omega} \hat{T} \{ \phi(x_1) \cdots \phi(x_N) \} \ket{\Omega},
\]
which is the Green function of the theory, we can compute how it scales under the transformation of the fields. We have:
\begin{enumerate}[label=(\alph*)]
    \item Since...
    \item For the amputated Green functions, we have...
    \item For the scattering matrix elements, we recall the LSZ reduction formula as...
\end{enumerate}